\section{Background and Related Work}
\label{sec:background}

\subsection{Digital twins: definitions and maturity}

\citet{grieves2017digital} defined the digital twin as a virtual representation linked to a physical product across its life cycle to expose emergent behaviour before deployment; industrial surveys broadened the concept to processes and services~\citep{tao2019}. \citet{jones2020characterising} consolidated thirteen twin characteristics and highlighted \emph{fidelity} as a central open gap.

Two distinctions govern our positioning. First, \citet{kritzinger2018digital} classify systems by data-flow automation: \emph{digital models} (manual both ways), \emph{digital shadows} (automatic physical-to-virtual only) and \emph{digital twins} (automatic both ways). Second, \citet{wright2020tell} argue that a twin requires an existing physical counterpart. For a pre-deployment architectural twin---where the production system does not yet exist or cannot safely endure destructive fault injection---we adopt Kritzinger's ladder explicitly (\S\ref{sec:maturity}) and answer Wright and Davidson by generating an \emph{executable real-middleware emulation} (real Cyclone DDS over Linux containers, distinct from discrete-event simulation) while explicitly bounding what that substitution can and cannot prove (\S\ref{sec:fidelity}). ISO~23247-1 defines manufacturing-oriented digital-twin requirements (accuracy, data acquisition, communication)~\citep{iso23247}; our requirements in \S\ref{sec:requirements} address analogous concerns for software architectures.

\subsection{Model-driven engineering of digital twins}

Model-driven engineering treats models as primary artefacts~\citep{schmidt2006mde}. \citet{bordeleau2020towards} identify heterogeneous model management, bidirectional synchronisation and life-cycle collaboration as core challenges for model-driven twins, while \citet{eramo2022conceptualizing} conceptualise coupling models with evolving operational data. \citet{kirchhof2020model} synthesise a twin's \emph{information system} from MontiArc and class diagrams; TwinOps~\citep{hugues2020twinops} combines AADL, Modelica and DevOps for CPS. \citet{dalibor2022cross} surveyed 356 publications and documented a scarcity of systematic engineering methods and rigorous evaluations. Unlike \citet{kirchhof2020model}, we generate the twin's \emph{physical stand-in} (an executable deployment) from the architecture model, making bidirectional synchronisation concrete via auditable M2S (\S\ref{sec:m2s}) and S2M (\S\ref{sec:s2m}) transformations.

\subsection{Runtime models and architectural analysis}

Models@run.time maintains a causally connected self-representation for live adaptation~\citep{blair2009models}, whereas our twin targets offline pre-deployment verification with telemetry-based calibration. Classical architecture-based dependability and performance approaches---reliability models~\citep{gosevapopstojanova2001architecture}, AADL Error-Model Annex~\citep{delange2014aadl}, Palladio~\citep{becker2009palladio} and static ROS recovery in ROSDiscover~\citep{timperley2022rosdiscover}---provide the model side but rarely generate an executable middleware counterpart or quantify fidelity against it. Closest to our S2M path, \citet{mazkatli2020incremental} incrementally calibrate Palladio parameters from monitoring data and \citet{walter2017expandable} extract performance models from logs. Because our counterpart is generated from the model itself, our calibrator must additionally establish parameter identifiability (\S\ref{sec:identifiability}). Our metamodel extends the graph-based dependency derivation of \citet{yigit2025graph}.

\subsection{Emulation, chaos engineering and simulation validation}

Container- and namespace-based emulators such as Mininet~\citep{lantz2010network} and Containernet~\citep{peuster2016medicine} run distributed topologies on one machine. Chaos engineering~\citep{basiri2016chaos} and service-level fault injection such as Filibuster~\citep{meiklejohn2021filibuster} test resilience in running systems. We combine model-generated emulation with a shared fault vocabulary to create an observed reference for simulators. Finally, simulation-validation theory~\citep{sargent2013verification,kleijnen2015design} warns against circular validation against shared assumptions, and computer-model calibration~\citep{kennedy2001bayesian} stresses explicit uncertainty and model discrepancy. We transfer both principles to pre-deployment twins via an explicit independence analysis and pre-registered claim table (\S\ref{sec:fidelity}).
